%!TEX root = main.tex
\vspace{-3.7mm}
\begin{abstract}
Semidefinite Programming (SDP) and Sums-of-Squ- ares (SOS) relaxations have 
led to
%allowed the design of 
\emph{certifiably} optimal \emph{non-minimal} solvers for several robotics and computer vision problems. 
However, most non-minimal solvers rely on least squares formulations, and, as a result, are 
brittle against
%not robust to
outliers.
%, 
%producing poor results in practice.
% , and limiting their applicability to the outlier-free case.
%
While a standard approach to regain robustness against outliers 
is to use robust cost functions, the latter 
% To overcome this challenge, 
% %On the other hand, 
% the robust estimation literature
% has developed
% %provides 
% a grounded theory to gain robustness against outliers,
% reformulating the least squares problems
% %through
% with
%  the use of robust cost functions. 
% But robust cost functions 
typically introduce 
%extra
other non-convexities, 
preventing 
the use of 
existing 
%the
non-minimal solvers.
 % 
In this paper, we 
%bridge the gap between
enable the simultaneous use of 
non-minimal solvers and robust estimation
%
by providing a general-purpose approach for robust \emph{global} estimation, which can be applied to any {problem where a non-minimal solver is available for the outlier-free case.}
% is available.
%Towards this goal, 
To this end,
we leverage the Black-Rangarajan duality between robust estimation and outlier processes (which has been traditionally applied to early vision problems), and show that
%that 
%the 
 % joint use of non-minimal solvers and \emph{graduated non-convexity} (\GNC) 
%use of 
\emph{graduated non-convexity} (\GNC) 
can be used 
in conjunction with non-minimal solvers 
%allows 
%computing 
to compute robust solutions, without requiring 
an initial guess. % for the variables. 
 %
\revise{Although \GNC's global optimality cannot be guaranteed, }we demonstrate\revise{~the empirical robustness of} the resulting \emph{robust non-minimal solvers} in applications, including 
point cloud and mesh registration, pose graph optimization, and image-based object pose estimation (also called \emph{shape alignment}).
 % Replacing the least squares cost with robust cost functions yields non-convex optimizations whose performance are highly sensitive to initial guesses. 
Our solvers are robust to {70-80\%} of outliers, 
%across applications, 
outperform \ransac, 
are more accurate than specialized \emph{local} solvers, and 
faster than specialized \emph{global} solvers.
% and is more broadly applicable 
% than existing general-purpose methods (\eg{ }\ransac), and compares favorably against specialized methods, while
% % still, it is
% being 
% significantly
% faster.
% and general-purpose. 
\finalize{We also propose the \emph{first} certifiably optimal non-minimal solver for shape alignment using SOS relaxation.}
% We also extend the literature on non-minimal solvers by proposing a certifiably optimal SOS solver for shape alignment.
% image-based object pose estimation (also called \emph{shape alignment}).
% faster than specialized methods
% more robust then RANSAC and existing techniques
 % The third contribution is to demonstrate how two examples of the \GNC framework, the \GNC-\emph{Geman McClure} (\GNCGM) and the \GNC-\emph{Truncated Least Squares} (\GNCTLS), can be applied on four robotics and computer vision applications: point cloud registration, generalized 3D registration, pose graph optimization, and shape alignment from a single image. 
 %
 % Extensive benchmarks on both simulated and real datasets show the \GNC framework achieves state-of-the-art robustness against 70\% to 80\% outliers across applications, while being a general framework. Lastly, our development of a globally optimal solver based on SOS relaxation for the single shape estimation from a single image is itself of interest.
\end{abstract}

% Keywords appear just beneath the abstract. Use only for final RAL version.
\begin{IEEEkeywords}Graduated non-convexity, outlier rejection, robust estimation, spatial perception, global optimization.
\end{IEEEkeywords}